% !PRINTED 399
\weberpage{399}
%<<CONT
(i.e. of the coefficients in its equation) there follows this one, that
the system which one obtains, when every imaginary line is
replaced by its conjugate, is again an Aronhold system
of the same curve. For then the coefficients do not change
when $i$ is replaced by $-i$.

The curve therefore certainly becomes real if, in the given
system of seven, the conjugate of every imaginary line
occurs; then they form a \gesperrt{real system of seven}.

Such a real system can now contain:

1) one real and three pairs of conjugate imaginary lines;

2) three real and two pairs of conjugate imaginary lines;

3) five real and one pair of conjugate imaginary lines;

4) seven real lines.

We shall see that these four assumptions lead, in the same
order, to the four cases enumerated in the previous
section.

This demonstration becomes very simple if one makes use of the
mode of designating the double tangents which we set forth in \S.\,99,
according to which we can form the Steiner complexes
immediately from the designation.

1) In the first case we designate the given seven
lines by
\[
0,\ 1,\ 1',\ 2,\ 2',\ 3,\ 3',
\]
assume $0$ to be real, and $1$ conjugate imaginary with $1'$, $2$ with $2'$,
$3$ with $3'$. The remaining 21 double tangents are then designated
by the symbols $[01], [01'], [11'], \dots$.

According to the rule of \S.\,99 we obtain the two
following Steiner complexes:
\[
\begin{aligned}
(1)&\quad 0[01],\ 1'[11'],\ 2[12],\ 2'[12'],\ 3[13],\ 3'[13']\\
(1')&\quad 0[01'],\ 1[11'],\ 2[1'2],\ 2'[1'2'],\ 3[1'3],\ 3'[1'3'].
\end{aligned}
\]

The complex conjugate to (1) must contain the elements $0,1,2,2',3,3'$,
and is therefore, by \S.\,98, 7., identical with (1'), which
likewise contains these six elements.

From this it follows that $[11']$ is real, and that
\[
\begin{aligned}
&{}[01]\quad [01']\\
&{}[12]\quad [1'2']\\
&{}[12']\quad [1'2]
\end{aligned}
\]
%CONT>>
